\section{Residuals, approximate certificates, and reactive stability}
\label{sec:numerical}

Exact feasibility and small residuals answer different questions. We make
their relationship explicit, first for the arithmetic reduction and then
for general bounded resistive networks. Throughout this section the voltage
bounds remain exact, including singleton bounds at pinned buses. Only
injection constraints are relaxed.

For a nonempty rational voltage box $B=\prod_i[\ell_i,u_i]$, define
\begin{equation}
 \rho_G(v)=\max_i\operatorname{dist}
       (\power_i(v),[p_i^-,p_i^+]),\qquad
 \gamma_G=\min_{v\in B}\rho_G(v).
 \label{eq:network-residual}
\end{equation}
Continuity and compactness imply that the minimum is attained and that
$\gamma_G=0$ exactly when the network is feasible. For a normalized source
instance $\Phi$, define
\begin{equation}
 R_\Phi(x)=\max\bigl\{|x+y-z|:\ x+y=z\text{ in }\Phi;
             \ |xy-1|:\ xy=1\text{ in }\Phi\bigr\},
 \qquad \gamma_\Phi=\min_{x\in[1/2,2]^n}R_\Phi(x).
 \label{eq:source-residual}
\end{equation}
A maximum over no constraints is zero, including for an empty network.

\subsection{A quantitative correspondence for the actual gadgets}

\begin{theorem}
\label{thm:residual-transfer}
For the ordinary construction of Section~\ref{sec:resistive}, with $m$
normalized equations,
\begin{equation}
 \frac{\gamma_\Phi}{10(6m+1)}\ \le\ \gamma_G
       \ \le\ 2\gamma_\Phi.
 \label{eq:residual-transfer}
\end{equation}
\end{theorem}
\begin{proof}
Given any source-box point $x$, extend it using exact complementary copies,
$v_I=x$, and $v_W=h(x)=2x-1+1/x$, whether or not it solves $\Phi$.
The voltage bounds hold by~\eqref{eq:W-range}. The free-injection bounds
are redundant throughout the box. All copy, $C_I$, and $I$ equations hold
exactly. An addition bus has residual $|x+y-z|$, and an inversion bus $D$
has residual $|y-1/x|\le2|xy-1|$. Taking a minimizing source point gives
the upper bound.

Conversely, let a network-box point have residual $\varepsilon$ and read
each source variable at its root. A copy pin has voltage exactly $1$, so its
injection residual bounds the error in its neighboring complement sum by
$\varepsilon$. There are at most $6m$ path extensions in the whole network.
Induction along a path therefore shows that each requested copy differs from
its intended root value or complement by at most
$\delta=6m\varepsilon$. The addition equation now gives source residual
at most $\varepsilon+3\delta$.

In an inversion gadget, write $x,y$ for the two root values. The residual
at $C_I$ gives $|v_I-x|\le\varepsilon+\delta$. Dividing the residual at
$I$ by $v_I\ge1/2$ gives
\[
 |v_W-h(v_I)|\le2\varepsilon.
\]
On $[1/2,2]$, $|h'|\le2$, so
$|v_W-h(x)|\le2\varepsilon+2(\varepsilon+\delta)$.
The residual at $D$, including its two weighted $x$ copies and one $y$
copy, gives
\[
 |y-(v_W-2x+1)|\le\varepsilon+3\delta.
\]
Combining these inequalities and $x\le2$ yields
\[
 |xy-1|\le2\bigl(\varepsilon+3\delta+2\varepsilon
                        +2(\varepsilon+\delta)\bigr)
           =10\varepsilon+10\delta.
\]
This also bounds every addition residual. Apply the resulting inequality
$\gamma_\Phi\le10(6m+1)\rho_G(v)$ to a minimizing network-box point.
If $m=0$, both residual minima are zero and the conclusion still holds.
\end{proof}

The theorem bounds the loss in transferring a residual; it does not supply
a lower bound on $\gamma_\Phi$ for arbitrary infeasible formulas. In fact,
even the fixed-data networks of Section~\ref{sec:resistive} can have
extremely small positive residual minima.

\subsection{Infeasible fixed-data networks with doubly exponential accuracy}

\begin{theorem}
\label{thm:tiny-residual}
For every integer $k\ge0$ there is an infeasible connected network $G_k$
with the fixed numerical data and degree bound of
Theorem~\ref{thm:rpf-complete}, at most $68k+67$ buses and $72k+72$
lines, and a rational voltage vector in its exact voltage box for which
every injection constraint except one holds exactly and
\begin{equation}
 0<\gamma_{G_k}\le\rho_{G_k}(v)<2^{-2^k}.
 \label{eq:tiny-residual}
\end{equation}
The strict lower bound concerns the minimum; the exhibited residual is only
an upper bound on that minimum.
\end{theorem}
\begin{proof}
Use the variables $t,h,x_0,\ldots,x_k$ and
$a_j,b_j,c_j$ for $0\le j<k$, all in $[1/2,2]$. Impose
\begin{equation}
\begin{gathered}
 tt=1,\qquad h+h=t,\qquad t+h=x_0,\\
 a_j+h=x_j,\quad x_jb_j=1,\quad a_j+b_j=c_j,\quad x_{j+1}+h=c_j
                 \quad(0\le j<k),\\
 x_kt=1.
\end{gathered}
 \label{eq:tiny-system}
\end{equation}
The first three equations force $t=1$, $h=1/2$, $x_0=3/2$.
Each group of four then forces
\[
 x_{j+1}=x_j+1/x_j-1.
\]
Writing $x_j=1+\delta_j$ gives
$\delta_0=1/2$ and
$\delta_{j+1}=\delta_j^2/(1+\delta_j)>0$.
The last equation requires $x_k=1$, a contradiction.

Nevertheless, assign precisely these recurrence values and set
$a_j=x_j-1/2$, $b_j=1/x_j$, and $c_j=x_j+1/x_j-1/2$.
Then
\[
 1<x_j\le3/2,\quad 1/2<a_j\le1,\quad
 2/3\le b_j<1,\quad 3/2<c_j\le5/3,
\]
so every voltage-source bound holds. All source equations except the last
hold exactly. More explicitly, $\delta_j=1/d_j$, where
$d_0=2$ and $d_{j+1}=d_j(d_j+1)$. Induction gives
$d_k\ge2^{2^k}$.

Apply the ordinary electrical construction and its canonical extension to
this nonsolution. The proof of Theorem~\ref{thm:residual-transfer} shows
that only the final inversion bus $D$ violates its injection constraint,
by exactly
\[
 1-1/x_k=\frac{1}{d_k+1}<2^{-2^k}.
\]
All voltage intervals hold exactly. The source system is infeasible, hence
so is the network, and compactness of its voltage box makes its residual
minimum strictly positive. There are $4k+3$ variables and $4k+4$ equations;
the size bounds follow from Theorem~\ref{thm:rpf-complete}.
The source incidence graph is connected, and each variable path and equation
gadget is connected with all incidence attachments present. Thus the network
is connected as well.
\end{proof}

Its standard graph encoding has length $O((k+1)\log(k+2))$. Therefore no
bound of the form $2^{-\operatorname{poly}(s)}$, with $s$ the input length,
can be a universal lower bound on positive residual minima in this fixed-data
class. Ruling out these infeasible instances solely by a universal residual
threshold can require exponentially many accuracy bits in the number of
buses. This is not a lower bound on exact decision algorithms: the displayed
family has the short algebraic infeasibility proof just given.

There is also a general lower bound at the doubly exponential scale. We
record it to identify the scale of the preceding construction, using a
general theorem about polynomial minima.

\begin{proposition}
\label{prop:residual-separation}
For a rational network with $n\ge1$ buses, put $q=n+1$.
There is an explicitly constructible integer coefficient bound $H$ with
polynomial bit length in the input such that, with
$M=4n+2$, $\widetilde H=\max\{H,2q+2M\}$, and
$\tau=\lceil\log_2\widetilde H\rceil$, infeasibility implies
\begin{equation}
 \gamma_G\ge 2^{-q4^q(\tau+q+4)}.
 \label{eq:separation}
\end{equation}
For fixed-data bounded-degree networks this yields
$\log_2(1/\gamma_G)\le2^{O(n)}$ whenever $0<\gamma_G<1$.
Together with Theorem~\ref{thm:tiny-residual}, the worst-case scale of
this logarithm in that class is $2^{\Theta(n)}$.
\end{proposition}
\begin{proof}
Let $U=\max_i u_i$, $D=\max_i\sum_j g_{ij}$, and
$A=\max_i\max\{|p_i^-|,|p_i^+|\}$.
On the positive voltage box, $|\power_i(v)|\le U^2D$, so
$B_0=1+U^2D+A$ bounds $\rho_G$ from above. Consider the compact set
\[
 T=\{(v,t):v\in B,\ 0\le t\le B_0,\
             t\ge p_i^- -\power_i(v),\
             t\ge\power_i(v)-p_i^+\ (i\in N)\}.
\]
It is connected: each point joins vertically to $(v,B_0)$, and the whole
top face $B\times\{B_0\}$ lies in $T$ and is convex. The minimum of
the polynomial $t$ on $T$ is $\gamma_G$.
There are $M=4n+2$ weak polynomial inequalities, of degree at most two.
Clear denominators in each by a positive integer, and let $H\ge1$ bound
the absolute values of the resulting integer coefficients, also including
the objective $t$. Rational arithmetic gives $\log H$ polynomial in the
input length.

The polynomial-minimum bound of
\citet[Theorem 1.1]{JeronimoEtAl2013}, applied in dimension $q\ge2$
with even degree bound $d=2$, gives for a nonzero minimum
\[
 \gamma_G\ge
  \bigl(2^{4-q/2}\widetilde H\,2^q\bigr)^{-q4^q}.
\]
The logarithm of the base is at most $\tau+q+4$, proving
\eqref{eq:separation}. With a fixed rational data alphabet and bounded
degree, $B_0$ and the cleared coefficient bound can be chosen constant.
Then $\tau=O(\log(n+2))$, and the upper estimate on
$\log_2(1/\gamma_G)$ follows. For the lower estimate, the family above
has linearly many buses in $k$ (in particular its distinct variable roots
number $4k+3$), and its logarithmic residual exceeds $2^k$.
\end{proof}

\subsection{Short rational certificates for a gap promise}

\begin{proposition}
\label{prop:approx-certificates}
Let $\varepsilon>0$ be rational. Every exactly feasible rational $\RPF$
instance has a rational voltage vector in its exact box with
$\rho_G(v)\le\varepsilon/2$, of bit length polynomial in the network input
and the encoding length of $\varepsilon$. Consequently, the promise problem
whose yes cases have $\gamma_G=0$ and whose no cases have
$\gamma_G>\varepsilon$ has polynomial-size certificates verifiable in
deterministic polynomial time. No answer is prescribed when
$0<\gamma_G\le\varepsilon$.
\end{proposition}
\begin{proof}
An empty network is an immediate yes instance and needs no voltage
coordinates. Otherwise put $U=\max_i u_i$ and
$D=\max_i\sum_jg_{ij}$.
For $v,w$ in the voltage box, differentiating~\eqref{eq:rpf} gives
\[
 \left|\frac{\partial\power_i}{\partial v_i}\right|
       \le3UD,
 \qquad
 \sum_{j\ne i}\left|\frac{\partial\power_i}{\partial v_j}\right|
       \le UD.
\]
The segment between $v$ and $w$ remains in the box, so
\begin{equation}
 |\power_i(v)-\power_i(w)|\le4UD\|v-w\|_\infty.
 \label{eq:lipschitz}
\end{equation}
The distance to an interval is $1$-Lipschitz, and the same bound holds for
the change of $\rho_G$.
Choose a dyadic mesh $\eta=2^{-b}$ with
$\eta\le\varepsilon/(2\max\{1,4UD\})$ and $b\ge0$.
Round each coordinate of an exactly feasible point down to this mesh and
then clamp it to its own rational interval. The resulting point stays in
the box, moves by at most $\eta$, and has residual at most
$\varepsilon/2$. Singleton bounds are preserved by clamping. The exponent
$b$ and all resulting rational-coordinate lengths are polynomial in the
input encoding and in that of $\varepsilon$.

A verifier checks the exact rational voltage bounds and evaluates all
injections with rational arithmetic, accepting when the residual is at most
$\varepsilon/2$. The stated bit bound gives completeness on promised yes
instances. No certificate can be accepted when $\gamma_G>\varepsilon$,
giving soundness on promised no instances.
\end{proof}

This is a certificate-existence result, not an algorithm for finding such
a point. It does not put exact $\RPF$ in $\NP$, nor does it decide all
instances of an unpromised approximate feasibility problem at a sharp
tolerance. In particular, it does not resolve the approximate-membership
question of \citet[Section 1.3 of the arXiv version]{BienstockVerma2019}:
their question relaxes sine equations in a lossless, fixed-magnitude model
and does not impose the promise used here.

\subsection{Stability of the equal-angle transfer}

The exact reactive argument also has a quantitative version. It applies to
real bus angles with short edge differences; a principal-angle winding
solution does not satisfy that hypothesis.

\begin{proposition}
\label{prop:reactive-stability}
Consider a connected purely resistive AC network with no shunts and $n\ge2$
buses. Let $0<\ell\le v_i\le U$, and suppose the real angle differences on
all lines lie in $[-\pi/2,\pi/2]$. Center the angles so that
$\sum_i\theta_i=0$, let $q$ be their reactive-injection vector, and let
$\lambda_2>0$ be the smallest positive eigenvalue of the conductance
Laplacian $L_g$. Then
\begin{align}
 \|\theta\|_2&\le
       \frac{\pi\|q\|_2}{2\ell^2\lambda_2},
       \label{eq:angle-stability}\\
 0\le P_i(v,\theta)-\power_i(v)&\le
       \frac{\pi^2U^2\|q\|_2^2}{8\ell^4\lambda_2}
       \qquad(i\in N).
       \label{eq:active-stability}
\end{align}
One may replace the constants $\pi/2$ and $\pi^2/8$ by the rational
constants $2$ and $2$, respectively. Moreover,
\begin{equation}
 \lambda_2\ge\frac{2g_{\min}}{(n-1)^2},
 \qquad g_{\min}=\min_{\{i,j\}\in E}g_{ij}.
 \label{eq:spectral-lower}
\end{equation}
Thus entirely rational bounds can be computed directly from the graph data.
For disconnected networks these statements apply componentwise. An isolated
bus has $q_i=0$ and zero active discrepancy.
\end{proposition}
\begin{proof}
Put $s=\theta^TL_g\theta=\sum_{\{i,j\}}g_{ij}(\theta_i-\theta_j)^2$.
The energy identity~\eqref{eq:angle-energy} and
$t\sin t\ge(2/\pi)t^2$ for $|t|\le\pi/2$ give
\[
 (2/\pi)\ell^2s\le-\theta^Tq
       \le\|\theta\|_2\|q\|_2
       \le\sqrt{s/\lambda_2}\,\|q\|_2.
\]
If $s=0$, the centered angles vanish and both conclusions are immediate.
Otherwise divide by $\sqrt s$ to obtain
$s\le\pi^2\|q\|_2^2/(4\ell^4\lambda_2)$, which also gives
\eqref{eq:angle-stability}. From~\eqref{eq:resistive-ac-P} and
$0\le1-\cos t\le t^2/2$,
\[
 0\le P_i(v,\theta)-\power_i(v)
   =\sum_jg_{ij}v_iv_j(1-\cos(\theta_i-\theta_j))
   \le (U^2/2)s.
\]
This proves~\eqref{eq:active-stability}. Using the weaker inequality
$t\sin t\ge t^2/2$ gives the stated rational constants.

Finally, along a simple path of at most $n-1$ edges, Cauchy--Schwarz gives
$(\theta_i-\theta_j)^2\le(n-1)s/g_{\min}$.
For centered angles,
\[
 \|\theta\|_2^2=\frac1n\sum_{i<j}(\theta_i-\theta_j)^2
       \le\frac{(n-1)^2}{2g_{\min}}s.
\]
Minimizing the Rayleigh quotient over nonzero centered vectors proves
\eqref{eq:spectral-lower}. Components can be centered independently because
the electrical quantities depend only on edge differences.
\end{proof}

\begin{corollary}
\label{cor:ac-residual-transfer}
Consider an AC instance obtained from the ordinary $m$-equation resistive
construction, with $n\ge1$ buses, real line differences in
$[-\pi/2,\pi/2]$, and exact voltage bounds. If a state has active-injection
interval violation at most $\varepsilon$ and $\|q\|_\infty\le\eta$, its
root voltage assignment has source residual at most
\begin{equation}
 10(6m+1)\bigl(\varepsilon+256n(n-1)^2\eta^2\bigr).
 \label{eq:ac-residual-transfer}
\end{equation}
\end{corollary}
\begin{proof}
Here $\ell=1/2$, $U=4$, and $g_{\min}\ge1$ in every nonsingleton
component. The rational bounds of Proposition~\ref{prop:reactive-stability},
and $\|q\|_2^2\le n\eta^2$, bound the active discrepancy by
$256n(n-1)^2\eta^2$, also for smaller components. Subtracting this
discrepancy from the actual AC injections bounds $\rho_G(v)$ by the
parenthesis in~\eqref{eq:ac-residual-transfer}. The pointwise reverse
estimate proved in Theorem~\ref{thm:residual-transfer} gives the result.
Singleton components have zero discrepancy, including the case $n=1$.
\end{proof}

These estimates quantify the equal-angle argument by elementary energy and
Poincar\'e inequalities. They control the error induced by a small reactive
residual, but establish neither hardness with symmetric positive reactive
tolerances nor the absence of the winding solutions in
Proposition~\ref{prop:winding-counterexample}.
